\chapter{Уравнения, системы и неравенства}
\begin{definition}\term{Корень уравнения} (\emph{root; solution of an equation}) — допустимое значение неизвестного, при котором равенство верно.\end{definition}
\begin{definition}\term{Решение неравенства} (\emph{solution of an inequality}) — значение неизвестного, при котором указанное сравнение верно.\end{definition}
\begin{definition}\term{Решение системы} (\emph{solution of a system}) — значения неизвестных, которые удовлетворяют всем её условиям одновременно.\end{definition}
\section*{Линейные и квадратные уравнения}
\(ax + b = 0\) при \(a \ne  0\) имеет корень \(x = - \frac{b}{a}\). При \(a = 0,\allowbreak b = 0\) подходят все \(x\); при \(a = 0,\allowbreak b \ne  0\) решений нет.\par
\begin{definition}\term{Дискриминант} (\emph{discriminant}) квадратного уравнения \(ax^{2}+bx+c=0\) при \(a\ne0\) — число \(D=b^{2}-4ac\).\end{definition}
При \(D > 0\) два корня: \(x = \frac{- b \pm  \sqrt{D}}{2 a}\). При \(D = 0\) один корень: \(- \frac{b}{2 a}\). При \(D < 0\) действительных корней нет.\par
\section*{Правила и проверка}
К обеим частям равенства или неравенства можно прибавить одно и то же выражение при сохранении ОДЗ.\par
При умножении или делении неравенства на отрицательное число знак меняется; на положительное — сохраняется.\par
Метод интервалов: отметь нули и запрещённые точки, определи знак на каждом участке и отдельно проверь включение границ.\par
При возведении уравнения в квадрат могут появиться посторонние корни: проверь найденные значения в исходном условии.\par
{\small\textit{Пример: \(- 2 x > 6\) даёт \(x < - 3\). Уравнение \(\sqrt{x + 2} = x\) требует \(x \ge  0\) ещё до возведения в квадрат.}}\par

% Дополнение к главе 15. Только действительные решения; формулы Кардано и Феррари не используются.
\subsection*{Многочлены и уравнения третьей степени}
\begin{definition}
\term{Многочлен} (polynomial) от \(x\) — выражение
\[
P(x)=a_nx^n+a_{n-1}x^{n-1}+\cdots+a_1x+a_0.
\]
Если \(a_n\neq0\), число \(n\) называют \term{степенью многочлена} (degree). \term{Корень многочлена} (polynomial root) — число \(r\), для которого \(P(r)=0\).
\end{definition}

\begin{definition}
Корень \(r\) имеет \term{кратность} (multiplicity) \(k\), если \(P(x)=(x-r)^kQ(x)\), где \(Q(r)\neq0\). При \(k=1\) корень называют простым, при \(k>1\) — кратным.
\end{definition}

\begin{theorem}[Сумма и разность кубов]
Для любых действительных \(a,b\)
\[
a^3-b^3=(a-b)(a^2+ab+b^2),\qquad
a^3+b^3=(a+b)(a^2-ab+b^2).
\]
\end{theorem}
\begin{proof}
При раскрытии скобок в первом произведении получаем
\(a^3+a^2b+ab^2-a^2b-ab^2-b^3=a^3-b^3\).
Во втором произведении средние слагаемые также сокращаются:
\(a^3-a^2b+ab^2+a^2b-ab^2+b^3=a^3+b^3\).
\end{proof}

\begin{theorem}[Теорема Безу]
\term{Теорема Безу} (factor theorem): для любого действительного \(r\) существует многочлен \(Q(x)\), для которого
\[
P(x)=(x-r)Q(x)+P(r).
\]
Следовательно, \(r\) — корень \(P\) тогда и только тогда, когда \(P(x)\) делится на \(x-r\) без остатка.
\end{theorem}
\begin{proof}
Для натурального \(k\)
\[
x^k-r^k=(x-r)(x^{k-1}+x^{k-2}r+\cdots+r^{k-1}).
\]
Равенство проверяется умножением: все промежуточные слагаемые сокращаются. Поэтому каждое слагаемое суммы \(P(x)-P(r)=\sum_{k=1}^n a_k(x^k-r^k)\) делится на \(x-r\). Это даёт нужный \(Q\). Если \(P(r)=0\), остаток равен нулю. Обратно, подставив \(x=r\) в \(P(x)=(x-r)Q(x)\), получаем \(P(r)=0\).
\end{proof}

\begin{theorem}[Рациональный корень]
Пусть коэффициенты \(P(x)=a_nx^n+\cdots+a_0\) целые, \(a_n\neq0\), \(a_0\neq0\). Если \(\frac pq\) — \term{рациональный корень} (rational root), причём \(p,q\) взаимно просты и \(q>0\), то \(p\) делит \(a_0\), а \(q\) делит \(a_n\).
\end{theorem}
\begin{proof}
Подстановка \(\frac pq\) и умножение на \(q^n\) дают
\[
a_np^n+a_{n-1}p^{n-1}q+\cdots+a_1pq^{n-1}+a_0q^n=0.
\]
Все слагаемые, кроме последнего, делятся на \(p\), поэтому \(p\) делит \(a_0q^n\). Поскольку \(p\) и \(q^n\) взаимно просты, \(p\) делит \(a_0\). Аналогично \(q\) делит \(a_np^n\), а значит, делит \(a_n\). Использованное свойство взаимной простоты следует из разложения на простые множители: множители делителя, отсутствующие в одном сомножителе, должны целиком входить в другой.
\end{proof}

Если старший коэффициент равен \(1\), рациональный корень может быть только целым делителем свободного члена. Проверяем кандидатов подстановкой, а найденный множитель \(x-r\) отделяем по теореме Безу. Если свободный член нулевой, сначала выносим \(x\). Отсутствие рациональных корней не означает отсутствия действительных: например, \(x^3-2=0\) имеет корень \(\sqrt[3]2\).

\begin{example}
Для \(P(x)=x^3-4x^2+x+6\) проверка даёт \(P(2)=0\). Поэтому
\[
P(x)=(x-2)(x^2-2x-3)=(x-2)(x-3)(x+1).
\]
Уравнение \(P(x)=0\) имеет корни \(-1,2,3\).
\end{example}

\begin{theorem}[Устранение квадратного слагаемого]
В \term{кубическом уравнении} (cubic equation)
\(ax^3+bx^2+cx+d=0\), где \(a\neq0\), замена \(x=t-\frac{b}{3a}\) даёт \term{неполное кубическое уравнение} (depressed cubic)
\[
t^3+pt+q=0,\qquad
p=\frac{3ac-b^2}{3a^2},\qquad
q=\frac{2b^3-9abc+27a^2d}{27a^3}.
\]
\end{theorem}
\begin{proof}
После деления на \(a\) обозначим коэффициенты через \(B=\frac ba\), \(C=\frac ca\), \(D=\frac da\). Подставим \(x=t-\frac B3\) и раскроем скобки:
\[
\begin{aligned}
&\left(t-\frac B3\right)^3+B\left(t-\frac B3\right)^2+C\left(t-\frac B3\right)+D\\
&\qquad=t^3+\left(C-\frac{B^2}{3}\right)t+\frac{2B^3}{27}-\frac{BC}{3}+D.
\end{aligned}
\]
Коэффициент при \(t^2\) равен нулю. Подстановка \(B,C,D\) даёт указанные \(p,q\).
\end{proof}

Такая замена упрощает форму уравнения; сама по себе она не даёт готовой формулы корней. В задачах этого раздела после замены достаточно разложения на множители или выделения куба.

\subsection*{Два вида уравнений четвёртой степени}
\begin{definition}
\term{Биквадратное уравнение} (biquadratic equation) имеет вид \(ax^4+bx^2+c=0\), где \(a\neq0\). Замена \(u=x^2\) приводит к квадратному уравнению \(au^2+bu+c=0\) с обязательным условием \(u\geq0\).
\end{definition}
Каждое \(u>0\) даёт два корня \(x=\pm\sqrt u\), значение \(u=0\) — один корень \(x=0\), а отрицательное \(u\) действительных корней не даёт.

\begin{example}
\(x^4-10x^2+9=0\) после замены \(u=x^2\) даёт \((u-1)(u-9)=0\). Поэтому \(x=-3,-1,1,3\).
\end{example}

\begin{definition}
\term{Возвратное уравнение четвёртой степени} (reciprocal quartic equation) в этом разделе — уравнение
\[
ax^4+bx^3+cx^2+bx+a=0,\qquad a\neq0,
\]
у которого коэффициенты симметричны относительно середины.
\end{definition}

\begin{theorem}[Замена для возвратного уравнения]
Корни указанного возвратного уравнения ненулевые. Замена \(t=x+\frac1x\) сводит поиск действительных корней к
\[
at^2+bt+c-2a=0,\qquad |t|\geq2.
\]
Для каждого допустимого \(t\) решают \(x^2-tx+1=0\).
\end{theorem}
\begin{proof}
При \(x=0\) левая часть исходного уравнения равна \(a\neq0\), поэтому деление на \(x^2\) допустимо. Получаем
\[
a\left(x^2+\frac1{x^2}\right)+b\left(x+\frac1x\right)+c=0.
\]
Из \(t=x+\frac1x\) следует \(x^2+\frac1{x^2}=t^2-2\), что даёт квадратное уравнение для \(t\). Обратная подстановка равносильна \(x^2-tx+1=0\). Его дискриминант \(t^2-4\) неотрицателен ровно при \(|t|\geq2\); свободный член \(1\) исключает нулевой корень. Значит, указанное условие необходимо и достаточно.
\end{proof}

\begin{example}
Для \(x^4-3x^3+2x^2-3x+1=0\) получаем \(t^2-3t=0\), то есть \(t=0\) или \(t=3\). Значение \(t=0\) не удовлетворяет \(|t|\geq2\). При \(t=3\) имеем \(x^2-3x+1=0\), откуда
\[
x=\frac{3-\sqrt5}{2}\quad\text{или}\quad x=\frac{3+\sqrt5}{2}.
\]
\end{example}

\Needspace{7\baselineskip}\section*{Задачи}
\begin{enumerate}[label={},leftmargin=3.2em,itemsep=10pt]
\Needspace{3\baselineskip}\item[\textbf{11.1.}] Реши \(3 x - 7 = 11\).
\Needspace{3\baselineskip}\item[\textbf{11.2.}] Реши \(2 (x - 3) = 5 x + 9\).
\Needspace{3\baselineskip}\item[\textbf{11.3.}] Реши \(x^{2} - 5 x + 6 = 0\).
\Needspace{3\baselineskip}\item[\textbf{11.4.}] Реши \(2 x^{2} - 8 = 0\).
\Needspace{3\baselineskip}\item[\textbf{11.5.}] Реши \(\frac{x + 1}{x - 2} = 2\).
\Needspace{3\baselineskip}\item[\textbf{11.6.}] Реши систему: \(x + y = 7;\allowbreak x - y = 1\).
\Needspace{3\baselineskip}\item[\textbf{11.7.}] Реши \(- 3 x > 12\).
\Needspace{3\baselineskip}\item[\textbf{11.8.}] Реши \(x^{2} - 9 \le  0\).
\Needspace{3\baselineskip}\item[\textbf{11.9.}] Реши \((x - 1) (x + 2) > 0\).
\Needspace{3\baselineskip}\item[\textbf{11.10.}] Реши \(\frac{x - 3}{x + 1} \le  0\).
\Needspace{3\baselineskip}\item[\textbf{11.11.}] Реши \(| 2 x - 1 | = 5\).
\Needspace{3\baselineskip}\item[\textbf{11.12.}] Реши \(\sqrt{x + 2} = x\).
\Needspace{3\baselineskip}\item[\textbf{11.13.}] Из жизни. Поездка на такси стоит \(150\) рублей за подачу и \(30\) рублей за километр. Счёт составил \(510\) рублей. Какое расстояние проехали?
\Needspace{3\baselineskip}\item[\textbf{11.14.}] Из жизни. За \(3\) одинаковых блокнота и ручку стоимостью \(80\) рублей заплатили \(440\) рублей. Сколько стоит блокнот?
\Needspace{3\baselineskip}\item[\textbf{11.15.}] Из жизни. Для покупки билетов по \(350\) рублей выделено \(2000\) рублей. Какое максимальное целое число билетов можно купить?
\Needspace{3\baselineskip}\item[\textbf{11.16.}\starmark] Реши \(\sqrt{2 x + 3} = x\).
\Needspace{3\baselineskip}\item[\textbf{11.17.}\starmark] Реши \(\frac{x - 2}{x + 3} \ge  1\).
\Needspace{3\baselineskip}\item[\textbf{11.18.}\starmark] У прямоугольного участка периметр \(34\) м, площадь \(60\) м\(^{2}\). Найди стороны.
\Needspace{3\baselineskip}\item[\textbf{11.19.}] Докажите формулы Виета: если \(x_1,x_2\) — действительные корни \(ax^2+bx+c=0\), \(a\neq0\), с учётом кратности, то \(x_1+x_2=-\frac ba\), \(x_1x_2=\frac ca\).
\Needspace{3\baselineskip}\item[\textbf{11.20.}] Докажите, что квадратное уравнение \(ax^2+bx+c=0\), \(a\neq0\), имеет действительный корень тогда и только тогда, когда \(D=b^2-4ac\geq0\). Используйте выделение полного квадрата.
\Needspace{3\baselineskip}\item[\textbf{11.21.}] Решите в действительных числах \(x^3-8=0\), разложив разность кубов на множители.
\Needspace{3\baselineskip}\item[\textbf{11.22.}] Решите в действительных числах \(x^3+27=0\), разложив сумму кубов на множители.
\Needspace{3\baselineskip}\item[\textbf{11.23.}] Решите \(x^3-6x^2+11x-6=0\). Найдите целый корень, затем выделите соответствующий линейный множитель.
\Needspace{3\baselineskip}\item[\textbf{11.24.}] Решите \(2x^3-3x^2-8x+12=0\) в действительных числах.
\Needspace{3\baselineskip}\item[\textbf{11.25.}\starmark] В уравнении \(x^3-6x^2+9x-4=0\) сделайте замену \(x=t+2\). Решите полученное уравнение и вернитесь к \(x\).
\Needspace{3\baselineskip}\item[\textbf{11.26.}] Решите биквадратное уравнение \(x^4-5x^2+4=0\).
\Needspace{3\baselineskip}\item[\textbf{11.27.}] Решите в действительных числах \(x^4+3x^2-4=0\), проверив допустимость найденных значений \(x^2\).
\Needspace{3\baselineskip}\item[\textbf{11.28.}\starmark] Решите возвратное уравнение \(x^4-5x^3+6x^2-5x+1=0\). Для каждого корня уравнения относительно \(t=x+\frac1x\) проверьте условие \(|t|\geq2\).
\end{enumerate}
